\documentclass{article}
\usepackage[margin=1in]{geometry}
\usepackage{mathtools, amsfonts, amsthm, xcolor, listings}

\title{HW 2}
\author{Sam Ly}

\begin{document}
\maketitle

\section*{Chapter 4}

\begin{enumerate}
    \item {
        The lexical grammars of Python and Haskell are not regular. What does 
        that mean, and why aren’t they?

        The lexical grammars of Python and Haskell are not regular because the 
        levels of indentation are significant. Regular languages can not have 
        ``n-length string of characters'' have significant semantic meaning.
    }
    \item {
        Aside from separating tokens—distinguishing print foo from printfoo—spaces 
        aren’t used for much in most languages. However, in a couple of dark 
        corners, a space does affect how code is parsed in CoffeeScript, Ruby, 
        and the C preprocessor. Where and what effect does it have in each of 
        those languages?

        In C, the macro system places significance on the spaces, because C 
        macros can either be a regular value (eg. number literal, string, etc.) 
        or a ``function'' that gets subsituted.

        \begin{lstlisting}
#define FOO(x) 2*x
#define BAR (x)
// Usage:
FOO(a); // changes to 2*a
BAR; // changes to (x), likely an error.
        \end{lstlisting}
    }
    \item {
        Our scanner here, like most, discards comments and whitespace since those 
        aren’t needed by the parser. Why might you want to write a scanner that 
        does not discard those? What would it be useful for?

        Many modern languages can actually make use of comments. For example, 
        JavaDoc can turn comments into documentation. The JavaDoc can also inform 
        the IDE's LSP, allowing for inline completions and context.
    }
    \item {
        Add support to Lox’s scanner for C-style /* ... */ block comments. Make 
        sure to handle newlines in them. Consider allowing them to nest. Is 
        adding support for nesting more work than you expected? Why?

        
    }
\end{enumerate}

\end{document}
